% Fichier généré automatiquement par tools/generate_corrections.py.
% Source : STAT/STAT-03-Demarche/95_ChasseNeige/95_ChasseNeige.tex
% Statut : partiel (1/2 question(s) renseignée(s), 0 reprise(s)).
\begin{corrigebox}[Corrigé partiel extrait de la banque]
\CorrigeQuestion{1}
\CorrigeACompleter{Aucun élément de réponse n'est présent dans la source d'origine pour cette question.}
\CorrigeQuestion{2}
\paragraph*{Graphe de liaisons}

On commence par faire les figures planes puis le graphe de liaisons.
\begin{enumerate}
\item On cherche les solides ou les ensembles de solides soumis à 2 glisseurs . Le problème étant plan, les pivots dont l'axe est perpendiculaire au plan sont des glisseurs. \{10+11\} est un ensemble soumis à 2 glisseurs.
\item On isole ensuite 7 et on rélise un théorème du moment statique en $H$ suivant $\vy{3}$.
\end{enumerate}

\begin{center}
\includegraphics[width=\linewidth]{cor_01}
\end{center}

\paragraph*{On isole le vérin \{10+11\}}
D'après le PFS, l'ensmble étant soumis à 2 glisseurs, on a donc $\torseurstat{T}{11}{7} = \torseurl{F\vz{11}}{\vect{0}}{J}$.

\paragraph*{On isole \{7\}}
BAME : 
\begin{itemize}
\item action de la neige;
\item action de la pesanteur;
\item action de la pièce 11;
\item action de la pièce 3;
\end{itemize}

\paragraph*{On réalise le TMS en $H$ en projection sur $\vy{3}$.}

$$
\vectm{H}{\text{neige}}{7} \cdot \vy{3} + \vectm{H}{\text{Pesanteur}}{7} \cdot \vy{3} + \vectm{H}{11}{7} \cdot \vy{3} + \underbrace{\vectm{H}{3}{7} \cdot \vy{3}}_{\vect{0}} =0
$$

$
\Rightarrow
\left(\vect{HQ}\wedge Q\vect{x_7} \right) \cdot \vy{3} 
+ \left(\vect{HG}\wedge -gP\vect{y_3} \right) \cdot \vy{3} 
+ \left(\vect{HJ}\wedge F\vect{z_{11}} \right)  \cdot \vy{3}=0
$

$
\Rightarrow
\left(\left(a\vx{3}+b\vy{3}+c\vz{3}\right)\wedge Q\vect{x_7} \right) \cdot \vy{3} 
+ \underbrace{\left( i\vz{7}\wedge -gP\vect{y_3} \right) \cdot \vy{3}}_{\vect{0}}
+ \left(h\vz{7}\wedge F\vect{z_{11}} \right)  \cdot \vy{3}=0
$

$
\Rightarrow
\left(\vy{3} \wedge \left(a\vx{3}+c\vz{3}\right)  \right) \cdot   Q\vect{x_7}
+ \left(h\vz{7}\wedge F\vect{z_{11}} \right)  \cdot \vy{3}=0
$
$
\Rightarrow
\left( -a\vz{3}+c\vx{3}\right)   \cdot   Q\vect{x_7}
+ hF \sin\left(\beta - \gamma\right)\vy{3}  \cdot \vy{3}=0
$
$
\Rightarrow
Q\left( a\sin \gamma+c\cos\gamma \right)  
+ hF \sin\left(\beta - \gamma\right)=0
$

Au final, $F= - \dfrac{Q\left( a\sin \gamma+c\cos\gamma \right)  }{ h \sin\left(\beta - \gamma\right)}$.

$F$ étant l'effort déployé par le vérin, et $S$ sa section, on a alors, $F=pS$ et 
$S= - \dfrac{Q\left( a\sin \gamma+c\cos\gamma \right)  }{p h \sin\left(\beta - \gamma\right)}$.
\end{corrigebox}
